%% kapa: exam=1 round=20 year=2009 period=1 form=A subject=civil_law subject_ko=민법 seq=27 number=27
%% source: Q-Net 감정평가사 기출문제 (cst00309) artlSeq=5024178 file=extracted/1008/5024178/2009년 제20회 감정평가사 1차시험 기출문제/1교시.hwp
%% conversion: kordoc 4.24.51 -> tools/gen1.py (auto); review=proofread-sonnet; flags=none
\begin{question}{27}
선량한 풍속 및 기타 사회질서에 반하여 무효인 경우가 \uline{아닌} 것은? (다툼이 있으면 판례에 의함)
\begin{choices}
  \item 다수의 세입자입주권을 투기의 목적으로 매수하는 행위
  \item 이중양도의 경우 제2의 양수인이 양도인의 배임행위에 적극 가담한 경우
  \item 부첩관계의 종료를 해제조건으로 하는 증여계약
  \item 사찰이 그 존립에 필요불가결한 재산인 임야를 증여하는 행위
  \item 변호사 아닌 자가 민사소송의 당사자를 승소시켜주고 그 대가로 소송물의 일부인 임야지분을 양수받기로 약정한 경우
\end{choices}
\end{question}
\answer{1}
